\section{The dimension of the formal power series ring}
\label{sec:dimension}

Let $A=T[[x]]$. Interchanging the $x$- and $c$-coefficients and
using~\eqref{eq:finite-basis} identifies $A$ with the subring
\begin{equation}\label{eq:array}
 A=\left\{\sum_{n\ge0}b_nc^n\in B[[c]]:
                         b_n\in E_{2^n-1}\text{ for all }n\right\}.
\end{equation}
Indeed, a series in $T[[x]]$ gives at each $c$-degree $n$ an
element of $V_n[[x]]=E_{2^n-1}$. Conversely, expanding each $b_n$
in $x$ gives coefficients in $V_n$, so each resulting
$x$-coefficient belongs to $T$. Multiplication agrees because
each coefficient of $x^i c^j$ involves only a finite sum.

\Needspace{5\baselineskip}
\begin{proposition}\label{prop:upper}
One has $\dim A\le r+2$.
\end{proposition}

\begin{proof}
Equation~\eqref{eq:array} gives $A\subseteq H$.
If $h=\sum b_nc^n\in H$ has bound $w(b_n)\le M2^n$, choose
$d$ with $2^d\ge M+1$. Then
$M2^n\le2^{n+d}-1$, so $c^dh\in A$. Consequently
\begin{equation}\label{eq:localization}
 A[1/c]=H[1/c],\qquad\dim A[1/c]\le r.
\end{equation}
Put $P=\m[[x]]$. Since $\m^3\subseteq cT$,
\[
 P^3\subseteq(cT)[[x]]=cA\subseteq P.
\]
The equality in the middle follows by dividing coefficients by
$c$. Since $A/P\cong k[[x]]$, we obtain
\begin{equation}\label{eq:closed-part}
 \sqrt{cA}=P,\qquad\dim(A/cA)=1.
\end{equation}
In a finite prime chain of $A$, the segment avoiding $c$ has
length at most $r$, and the segment containing $c$ has length at
most one. If both occur, there is one additional inclusion between
them. Thus every such chain has length at most $r+2$.
\end{proof}

\begin{lemma}\label{lem:exponentials}
There exist $f_1,\ldots,f_r\in xk[[x]]$ algebraically
independent over $k$.
\end{lemma}

\begin{proof}
Choose $\lambda_1,\ldots,\lambda_r\in k$ linearly independent
over $\mathbb Q$ and set
\[
 f_i(x)=\exp(\lambda_i x)-1
       =\sum_{n\ge1}\frac{\lambda_i^n}{n!}x^n.
\]
Such $\lambda_i$ exist: for $r>1$, a root $\alpha$ of the
Eisenstein polynomial $t^r-2$ gives $1,\alpha,\ldots,\alpha^{r-1}$;
for $r=1$ take $\lambda_1=1$.
A polynomial relation in the $f_i$, after an invertible polynomial
translation, would yield
\[
 \sum_{\nu}a_\nu\exp(\mu_\nu x)=0,
 \qquad \mu_\nu=\sum_i\nu_i\lambda_i,
\]
with finitely many distinct multi-indices $\nu$ and not all
$a_\nu\in k$ zero. The $\mu_\nu$ are distinct. List them as
$\mu_1,\ldots,\mu_s$. Formal differentiation at zero in orders
$0,\ldots,s-1$ gives a homogeneous linear system with determinant
\[
 \det(\mu_j^{i-1})_{1\le i,j\le s}
   =\prod_{1\le i<j\le s}(\mu_j-\mu_i)\ne0,
\]
a contradiction. All operations are formal in characteristic zero.
\end{proof}

\begin{proposition}\label{prop:lower}
The ring $A$ contains a prime chain of length $r+2$.
Thus $\dim T[[x]]=r+2$.
\end{proposition}

\begin{proof}
Choose the series in \cref{lem:exponentials} and put
\[
 \q_j=(u_1-f_1(x),\ldots,u_j-f_j(x))B\quad(1\le j\le r),
 \qquad\q_0=(0).
\]
Before localizing, substitution of $f_i(x)$ for $u_i$, $i\le j$,
defines a surjection
\[
 F[U]\longrightarrow F[u_{j+1},\ldots,u_r]
\]
with kernel $(u_1-f_1(x),\ldots,u_j-f_j(x))F[U]$.
Its kernel is disjoint from $S$: expanding a nonzero element of
$k[U]$ in the remaining variables would otherwise give a nonzero
polynomial relation over $k$ among $f_1,\ldots,f_j$.
Localization therefore gives proper prime ideals and a strict chain
\begin{equation}\label{eq:qchain}
 (0)=\q_0\subsetneq\q_1\subsetneq\cdots\subsetneq\q_r.
\end{equation}
For strictness, $u_j-f_j(x)$ is nonzero in the domain obtained
after the preceding substitution, and remains nonzero after
localization. Full substitution extends to $B\longrightarrow L$
with kernel $\q_r$, since every element of $S$ has nonzero image
in $F\subseteq L$. This map fixes $F$, so
\begin{equation}\label{eq:qF}
 \q_r\cap F=(0).
\end{equation}

The coefficientwise ideal $\q_j[[c]]\subseteq B[[c]]$ is prime,
since its quotient is $(B/\q_j)[[c]]$, a domain.
Contract these primes to $A$:
\[
 Q_j=\q_j[[c]]\cap A\qquad(0\le j\le r).
\]
They all avoid $c$, and $Q_0=(0)$.
Since $u_j=D_1u_j/D_1\in W_1$ and $f_j(x)\in F$, one has
\[
 c(u_j-f_j(x))\in Q_j\setminus Q_{j-1}
 \qquad(1\le j\le r),
\]
where membership in $A$ follows from~\eqref{eq:array}.
Thus the contracted chain is strict.

If $a\in Q_r$, its coefficient at $c^0$ belongs to $E_0=F$
by~\eqref{eq:array} and to $\q_r$ by definition. It is zero
by~\eqref{eq:qF}. This says exactly that every $x$-coefficient
of $a$ belongs to $\m$, so $Q_r\subseteq P=\m[[x]]$.
The inclusion is strict because $c\in P\setminus Q_r$.
Finally $P\subsetneq P+(x)$ is a strict prime inclusion, since
$A/P\cong k[[x]]$. Hence
\begin{equation}\label{eq:final-chain}
 (0)=Q_0\subsetneq Q_1\subsetneq\cdots\subsetneq Q_r
       \subsetneq P\subsetneq P+(x)
\end{equation}
has length $r+2$. Combining this with \cref{prop:upper} proves
the dimension equality.
\end{proof}

\begin{proof}[Proof of \cref{thm:local}]
For $m=2$ use the discrete valuation ring $k[[c]]$.
For $m\ge3$ set $r=m-2$ and use $T_r$ from~\eqref{eq:T}.
Its required dimensions follow from \cref{prop:local,prop:lower}.
\end{proof}
